\documentclass[12pt]{article}

\usepackage[a4paper,margin=1in]{geometry}
\usepackage{amsmath,amssymb}
\usepackage{newtxtext,newtxmath} % Times-style font
\usepackage{fancyhdr}
\usepackage{listings}
\usepackage{graphicx}
\usepackage{float}
\usepackage{algorithm}
\usepackage{algpseudocode}
\usepackage{xcolor}
\usepackage{booktabs}
\usepackage{array}

\setlength{\textfloatsep}{10pt}
\setlength{\floatsep}{10pt}

\setlength{\parindent}{0pt}
\setlength{\parskip}{5pt}

% Keep entire document at 12 pt
\AtBeginDocument{
    \fontsize{12pt}{14pt}\selectfont
}

\lstdefinestyle{code}{
    language=Python,
    basicstyle=\ttfamily\fontsize{10pt}{12pt}\selectfont,
    breaklines=true,
    breakatwhitespace=false,
    frame=single,
    numbers=left,
    numberstyle=\tiny\color{gray},
    keywordstyle=\color{blue},
    commentstyle=\color{gray},
    stringstyle=\color{green!50!black},
    columns=fullflexible,
    keepspaces=true,
    showstringspaces=false
}

\lstdefinestyle{output}{
    basicstyle=\ttfamily\fontsize{10pt}{12pt}\selectfont,
    breaklines=true,
    breakatwhitespace=false,
    frame=single,
    numbers=none,
    backgroundcolor=\color{gray!4},
    columns=fullflexible,
    keepspaces=true,
    showstringspaces=false
}

% ------------------------------------------------
% Header
% ------------------------------------------------

\pagestyle{fancy}
\fancyhf{}

\fancyhead[L]{Name: Kishan Sahu{\hspace{4cm}}}
\fancyhead[R]{Enrollment No.: P26DS017}

\renewcommand{\headrulewidth}{0.4pt}
\setlength{\headheight}{15pt}

\begin{document}

% ------------------------------------------------
% TITLE
% ------------------------------------------------

\begin{center}
\textbf{Computer Science and Engineering Department, SVNIT, Surat}\\
\textbf{M.Tech. I -- Semester I}\\
\textbf{Design and Analysis of Algorithm (CSDS103)}\\
\textbf{Lab Assignment: Comprehensive Algorithmic Analysis}\\
\textbf{Lab Assignment 10: Backtracking Algorithms}
\end{center}
\vspace{-4pt}
\hrule
\vspace{4pt}

\vspace{0.3cm}

\section{Introduction}
Backtracking is a systematic algorithmic technique for solving constraint satisfaction problems (CSPs) by incrementally building candidates to the solutions and abandoning a candidate (``pruning'' the branch) as soon as it is determined that the candidate cannot possibly lead to a valid full solution.

This laboratory assignment implements and evaluates backtracking algorithms across three classic problems with advanced constraints:
\begin{enumerate}
    \item \textbf{Experiment 1}: Constrained 8-Queens Problem (with forbidden cells and pre-placed queens)
    \item \textbf{Experiment 2}: Constrained Sum of Subsets Problem (with cardinality bounds, mandatory/excluded items, and branch pruning diagnostics)
    \item \textbf{Experiment 3}: List-Constrained Graph Coloring Problem (with vertex color lists, pre-colored vertices, and search statistics)
\end{enumerate}

All implementations are self-contained in Python and directly correspond to the accompanying interactive Jupyter Notebook \texttt{Backtracking\_Lab\_Assignment.ipynb}.

% ============================================================================
% EXPERIMENT 1
% ============================================================================
\section{Experiment 1: Constrained 8-Queens Problem (Question 1.1)}

\subsection{Problem Statement \& Requirements}
Design a backtracking algorithm to solve the 8-Queens problem on an $8 \times 8$ chessboard such that no two queens share the same row, column, or diagonal. The algorithm extends the standard problem by supporting:
\begin{itemize}
    \item \textbf{Forbidden cells}: A set of board coordinates $(r, c)$ where queens cannot be placed.
    \item \textbf{Pre-placed queens}: Fixed queen positions $(r, c)$ that must remain intact.
    \item \textbf{Input validation \& Contradiction rejection}: Contradictory inputs (e.g., pre-placed queens on forbidden cells or attacking each other) must be rejected immediately before search.
    \item \textbf{Immediate branch pruning}: The branch must be pruned as soon as a queen placement conflicts with any previous queen or forbidden cell.
    \item \textbf{Symmetry reduction}: Support reducing symmetric branches when the constraints permit.
    \item \textbf{Solution enumeration}: Print all valid board arrangements and report the total count.
\end{itemize}

\subsection{Algorithmic Formulation}
A board state is represented as a 1D array \texttt{board[0..7]}, where \texttt{board[r] = c} indicates that the queen in row $r$ is placed in column $c$. The backtracking search proceeds row-by-row from $r = 0$ to $7$:
\begin{itemize}
    \item If row $r$ has a pre-placed queen, we verify its safety with previously placed rows; if safe, we recurse to row $r+1$, else backtrack.
    \item If row $r$ is unconstrained, we test column candidates $c \in \{0, \dots, 7\}$ using the bounding function \texttt{is\_safe(r, c)}. A placement is safe if $(r, c) \notin \text{forbidden}$, $c \neq \text{board}[i]$, and $|r - i| \neq |c - \text{board}[i]|$ for all $0 \le i < r$.
\end{itemize}

\subsection{Python Source Code Implementation}
The complete implementation of the solver function and board display helper is shown below:

\begin{lstlisting}[style=code]
def solve_constrained_8_queens(forbidden=None, pre_placed=None, symmetry_reduction=False):
    """
    Backtracking solver for the Constrained 8-Queens Problem.
    
    Parameters:
    - forbidden: set of (row, col) tuples where queens cannot be placed.
    - pre_placed: dict {row: col} representing fixed queen positions.
    - symmetry_reduction: bool, if True and constraints allow, reduces symmetric branches.
    
    Returns:
    - solutions: list of board configurations (each is a list where index=row, value=col).
    """
    if forbidden is None:
        forbidden = set()
    if pre_placed is None:
        pre_placed = {}

    # Requirement 4: Reject contradictory inputs
    # Check if any pre-placed queen occupies a forbidden cell
    for r, c in pre_placed.items():
        if (r, c) in forbidden:
            raise ValueError(f"Contradictory input: Pre-placed queen at ({r}, {c}) is in a forbidden cell.")

    # Check if any two pre-placed queens attack each other
    placed_list = list(pre_placed.items())
    for i in range(len(placed_list)):
        for j in range(i + 1, len(placed_list)):
            r1, c1 = placed_list[i]
            r2, c2 = placed_list[j]
            if r1 == r2 or c1 == c2 or abs(r1 - r2) == abs(c1 - c2):
                raise ValueError(f"Contradictory input: Pre-placed queens ({r1}, {c1}) and ({r2}, {c2}) attack each other.")

    solutions = []
    board = [-1] * 8  # board[r] represents the column of the queen in row r

    def is_safe(r, c):
        # Forbidden position check
        if (r, c) in forbidden:
            return False
        # Requirement 5: Prune immediately when attacking a previously placed queen
        for prev_r in range(r):
            prev_c = board[prev_r]
            if prev_c == c or abs(prev_r - r) == abs(prev_c - c):
                return False
        return True

    def backtrack(r):
        # Base case: All 8 queens placed successfully
        if r == 8:
            solutions.append(list(board))
            return

        # If current row has a pre-placed queen, respect it
        if r in pre_placed:
            c = pre_placed[r]
            if is_safe(r, c):
                board[r] = c
                backtrack(r + 1)
                board[r] = -1
            return

        col_range = range(8)
        # Requirement 6: Symmetry reduction (reflect horizontally across middle columns if unconstrained)
        if r == 0 and symmetry_reduction and not pre_placed and not forbidden:
            col_range = range(4)  # Explore only first 4 columns of row 0

        for c in col_range:
            if is_safe(r, c):
                board[r] = c
                backtrack(r + 1)
                board[r] = -1

    backtrack(0)
    return solutions


def print_board(solution):
    """Prints an 8x8 chessboard arrangement cleanly."""
    for r in range(8):
        row_str = " ".join("Q" if solution[r] == c else "." for c in range(8))
        print(f"Row {r}: {row_str}")
    print()\end{lstlisting}

\subsection{Demonstration Test Code}
The demonstration code evaluates unconstrained 8-Queens, constrained boards with pre-placed queens and forbidden cells, and verifies rejection of contradictory input:

\begin{lstlisting}[style=code]
print("=== Experiment 1: Constrained 8-Queens Demonstrations ===")

# Test 1: Standard 8-Queens (no constraints)
all_solutions = solve_constrained_8_queens()
print(f"1. Standard 8-Queens total valid arrangements: {len(all_solutions)}")

# Test 2: Pre-placed Queen at (0, 0) and forbidden cells
forbidden = {(1, 2), (2, 4), (3, 1)}
pre_placed = {0: 0}  # Queen at row 0, col 0
constrained_solutions = solve_constrained_8_queens(forbidden=forbidden, pre_placed=pre_placed)
print(f"2. Constrained arrangements (pre-placed (0,0), forbidden={(1,2),(2,4),(3,1)}): {len(constrained_solutions)}")
print("\nFirst valid arrangement:")
print_board(constrained_solutions[0])

# Test 3: Contradictory Input Handling
print("3. Testing Contradictory Input Detection:")
try:
    # Two pre-placed queens in attack position: (0, 0) and (1, 1) on same diagonal
    solve_constrained_8_queens(pre_placed={0: 0, 1: 1})
except ValueError as e:
    print("  Successfully rejected:", e)

try:
    # Pre-placed queen on a forbidden cell
    solve_constrained_8_queens(forbidden={(0, 3)}, pre_placed={0: 3})
except ValueError as e:
    print("  Successfully rejected:", e)\end{lstlisting}

\subsection{Execution Results \& Output Logs}
The execution outputs obtained from the interactive notebook are as follows:

\begin{lstlisting}[style=output]
=== Experiment 1: Constrained 8-Queens Demonstrations ===
1. Standard 8-Queens total valid arrangements: 92
2. Constrained arrangements (pre-placed (0,0), forbidden=((1, 2), (2, 4), (3, 1))): 3

First valid arrangement:
Row 0: Q . . . . . . .
Row 1: . . . . Q . . .
Row 2: . . . . . . . Q
Row 3: . . . . . Q . .
Row 4: . . Q . . . . .
Row 5: . . . . . . Q .
Row 6: . Q . . . . . .
Row 7: . . . Q . . . .

3. Testing Contradictory Input Detection:
  Successfully rejected: Contradictory input: Pre-placed queens (0, 0) and (1, 1) attack each other.
  Successfully rejected: Contradictory input: Pre-placed queen at (0, 3) is in a forbidden cell.

\end{lstlisting}

\subsection{Analysis}
\begin{itemize}
    \item \textbf{Classical 8-Queens}: The algorithm discovers all 92 valid solutions on the unconstrained $8 \times 8$ board.
    \item \textbf{Constrained Search}: Placing a fixed queen at $(0,0)$ and restricting coordinates $\{(1,2), (2,4), (3,1)\}$ prunes the search space from 92 to exactly 3 valid arrangements, confirming immediate branch pruning.
    \item \textbf{Contradiction Rejection}: The algorithm rejects diagonal conflicts between pre-placed queens (e.g., $(0,0)$ and $(1,1)$) and forbidden cell overlaps before any recursive exploration.
\end{itemize}

% ============================================================================
% EXPERIMENT 2
% ============================================================================
\section{Experiment 2: Constrained Sum of Subsets Problem (Question 2.1)}

\subsection{Problem Statement \& Requirements}
Given a set of $n$ positive integers $S$ and a target sum $T$, find all subsets whose elements sum exactly to $T$. The algorithm extends classical subset sum backtracking with the following constraints:
\begin{itemize}
    \item \textbf{Cardinality constraints}: Minimum and maximum subset size bounds $[k_{\min}, k_{\max}]$.
    \item \textbf{Mandatory elements}: Specified element values that must appear in every valid subset.
    \item \textbf{Excluded elements}: Specified element values that cannot appear in any valid subset.
    \item \textbf{Positional index tracking}: Duplicate numerical values are distinguished by their indices.
    \item \textbf{Pruning criteria}:
    \begin{enumerate}
        \item Prune branch when current sum exceeds $T$.
        \item Prune branch when maximum possible sum of remaining eligible elements cannot reach $T$.
        \item Prune branch when cardinality bounds $[k_{\min}, k_{\max}]$ cannot be satisfied.
    \end{enumerate}
    \item \textbf{Benchmark input}:
    $S = \{3, 5, 6, 7, 9, 11, 12\}, \quad T = 21, \quad k_{\min} = 2, \quad k_{\max} = 4, \quad \text{Mandatory} = \{3\}, \quad \text{Excluded} = \{7\}$
    Every valid subset must be reported, and every pruned branch must be logged with its reason.
\end{itemize}

\subsection{Algorithmic Formulation}
The search tree is a binary decision tree where each node corresponds to the decision of including or excluding item $i \in S$. The state is given by $(\text{idx}, \text{current\_indices}, \text{current\_sum})$.
Before exploring subtrees:
\begin{enumerate}
    \item If $\text{current\_sum} = T$, we verify $k_{\min} \le |\text{current\_indices}| \le k_{\max}$ and check mandatory containment.
    \item If $\text{current\_sum} > T$, prune (monotonicity of positive integers).
    \item If $|\text{current\_indices}| \ge k_{\max}$, prune further expansion.
    \item If $|\text{current\_indices}| + \text{remaining\_elements} < k_{\min}$, prune (cardinality underflow).
    \item If $\text{current\_sum} + \sum_{j=\text{idx}}^{n-1} S[j] < T$, prune (insufficient remaining sum).
    \item If prospective sum $\text{current\_sum} + S[\text{idx}] > T$, prune the inclusion branch.
\end{enumerate}

\subsection{Python Source Code Implementation}
The complete implementation of the constrained sum of subsets solver is shown below:

\begin{lstlisting}[style=code]
def solve_constrained_sum_of_subsets(S, T, k_min, k_max, mandatory=None, excluded=None):
    """
    Backtracking solver for the Constrained Sum of Subsets Problem.
    
    Parameters:
    - S: list of positive integers.
    - T: target sum.
    - k_min: minimum permitted subset size.
    - k_max: maximum permitted subset size.
    - mandatory: list of element values that must appear.
    - excluded: list of element values that cannot appear.
    
    Returns:
    - valid_subsets: list of valid subsets (each as a list of values).
    - pruning_logs: list of (current_subset, reason) explaining pruned branches.
    """
    if mandatory is None:
        mandatory = []
    if excluded is None:
        excluded = []

    mand_set = set(mandatory)
    excl_set = set(excluded)

    # Requirement 4: Reject invalid input
    if mand_set & excl_set:
        raise ValueError("Invalid input: Mandatory and excluded elements overlap.")
    if k_min > k_max or k_min < 0:
        raise ValueError("Invalid input: k_min cannot exceed k_max.")
    for m in mand_set:
        if m not in S:
            raise ValueError(f"Invalid input: Mandatory element {m} is not in input set S.")

    # Requirement 3: Distinguish duplicate values by tracking (index, value)
    items = [(i, val) for i, val in enumerate(S) if val not in excl_set]

    valid_subsets = []
    pruning_logs = []

    def backtrack(idx, current_indices, current_sum):
        current_vals = [S[i] for i in current_indices]

        # Check if target sum is met
        if current_sum == T:
            # Check cardinality constraint
            if k_min <= len(current_indices) <= k_max:
                # Check mandatory elements requirement
                if mand_set.issubset(set(current_vals)):
                    valid_subsets.append(list(current_vals))
            # Since all elements are positive, adding further elements will exceed T
            return

        if idx >= len(items):
            return

        # Pruning 1: Current sum exceeds T
        if current_sum > T:
            pruning_logs.append((current_vals, f"Current sum ({current_sum}) exceeds target ({T})"))
            return

        # Pruning 7a: Cardinality limit exceeded
        if len(current_indices) >= k_max:
            pruning_logs.append((current_vals, f"Subset size reached maximum permitted cardinality k_max={k_max}"))
            return

        # Pruning 7b: Cannot satisfy minimum cardinality k_min
        remaining_count = len(items) - idx
        if len(current_indices) + remaining_count < k_min:
            pruning_logs.append((current_vals, f"Remaining elements ({remaining_count}) cannot reach minimum cardinality k_min={k_min}"))
            return

        # Pruning 6: Maximum possible sum of remaining eligible elements cannot reach T
        remaining_sum = sum(val for _, val in items[idx:])
        if current_sum + remaining_sum < T:
            pruning_logs.append((current_vals, f"Max achievable sum ({current_sum + remaining_sum}) is less than target ({T})"))
            return

        item_idx, item_val = items[idx]

        # Branch 1: Include items[idx]
        # Pruning 5: Prune immediately if adding item_val exceeds T
        if current_sum + item_val <= T:
            backtrack(idx + 1, current_indices + [item_idx], current_sum + item_val)
        else:
            pruning_logs.append((current_vals + [item_val], f"Adding element {item_val} gives sum {current_sum + item_val} > {T}"))

        # Branch 2: Exclude items[idx]
        backtrack(idx + 1, current_indices, current_sum)

    backtrack(0, [], 0)
    return valid_subsets, pruning_logs\end{lstlisting}

\subsection{Benchmark Test Execution Code}
The code executing the required test case $S = \{3, 5, 6, 7, 9, 11, 12\}, T = 21$ is given below:

\begin{lstlisting}[style=code]
print("=== Experiment 2: Benchmark Test Execution ===")

S = [3, 5, 6, 7, 9, 11, 12]
T = 21
k_min = 2
k_max = 4
mandatory = [3]
excluded = [7]

subsets, logs = solve_constrained_sum_of_subsets(S, T, k_min, k_max, mandatory, excluded)

print(f"Given Set S = {S}")
print(f"Target T = {T}")
print(f"Cardinality range: [{k_min}, {k_max}]")
print(f"Mandatory elements: {mandatory}")
print(f"Excluded elements: {excluded}")
print("-" * 50)

print(f"Total Valid Subsets Found: {len(subsets)}")
for idx, s in enumerate(subsets, 1):
    print(f"  Valid Subset {idx}: {s} (Sum = {sum(s)}, Size = {len(s)})")

print("\n" + "-" * 50)
print(f"Pruned Branches Report (Total Pruned: {len(logs)}):")
print("-" * 50)
for branch, reason in logs:
    branch_str = str(branch) if branch else "[]"
    print(f"Branch {branch_str.ljust(22)} -> {reason}")\end{lstlisting}

\subsection{Execution Results \& Output Logs}
The complete execution results, valid subsets, and detailed pruning logs are shown below:

\begin{lstlisting}[style=output]
=== Experiment 2: Benchmark Test Execution ===
Given Set S = [3, 5, 6, 7, 9, 11, 12]
Target T = 21
Cardinality range: [2, 4]
Mandatory elements: [3]
Excluded elements: [7]
--------------------------------------------------
Total Valid Subsets Found: 1
  Valid Subset 1: [3, 6, 12] (Sum = 21, Size = 3)

--------------------------------------------------
Pruned Branches Report (Total Pruned: 29):
--------------------------------------------------
Branch [3, 5, 6, 9]           -> Adding element 9 gives sum 23 > 21
Branch [3, 5, 6, 11]          -> Adding element 11 gives sum 25 > 21
Branch [3, 5, 6, 12]          -> Adding element 12 gives sum 26 > 21
Branch [3, 5, 9, 11]          -> Adding element 11 gives sum 28 > 21
Branch [3, 5, 9, 12]          -> Adding element 12 gives sum 29 > 21
Branch [3, 5, 11, 12]         -> Adding element 12 gives sum 31 > 21
Branch [3, 5]                 -> Max achievable sum (20) is less than target (21)
Branch [3, 6, 9, 11]          -> Adding element 11 gives sum 29 > 21
Branch [3, 6, 9, 12]          -> Adding element 12 gives sum 30 > 21
Branch [3, 6, 11, 12]         -> Adding element 12 gives sum 32 > 21
Branch [3, 9, 11]             -> Adding element 11 gives sum 23 > 21
Branch [3, 9, 12]             -> Adding element 12 gives sum 24 > 21
Branch [3, 11, 12]            -> Adding element 12 gives sum 26 > 21
Branch [3]                    -> Max achievable sum (15) is less than target (21)
Branch [5, 6, 9, 11]          -> Adding element 11 gives sum 31 > 21
Branch [5, 6, 9, 12]          -> Adding element 12 gives sum 32 > 21
Branch [5, 6, 11]             -> Adding element 11 gives sum 22 > 21
Branch [5, 6, 12]             -> Adding element 12 gives sum 23 > 21
Branch [5, 9, 11]             -> Adding element 11 gives sum 25 > 21
Branch [5, 9, 12]             -> Adding element 12 gives sum 26 > 21
Branch [5, 11, 12]            -> Adding element 12 gives sum 28 > 21
Branch [5]                    -> Max achievable sum (17) is less than target (21)
Branch [6, 9, 11]             -> Adding element 11 gives sum 26 > 21
Branch [6, 9, 12]             -> Adding element 12 gives sum 27 > 21
Branch [6, 11, 12]            -> Adding element 12 gives sum 29 > 21
Branch [6]                    -> Max achievable sum (18) is less than target (21)
Branch [9, 11, 12]            -> Adding element 12 gives sum 32 > 21
Branch [11, 12]               -> Adding element 12 gives sum 23 > 21
Branch []                     -> Remaining elements (1) cannot reach minimum cardinality k_min=2

\end{lstlisting}

\subsection{Pruning Diagnostics Analysis}
On the test set $S=\{3,5,6,7,9,11,12\}$ with excluded element $7$ and mandatory element $3$:
\begin{itemize}
    \item \textbf{Unique Valid Solution}: Exactly one subset meets all constraints:
    $$\{3, 6, 12\} \quad (\text{Sum} = 21, \quad \text{Size} = 3 \in [2, 4], \quad 3 \in \text{subset}, \quad 7 \notin \text{subset})$$
    \item \textbf{Sum Exceeded Pruning}: Branches such as $[3, 5, 6, 9]$ were pruned because adding $9$ produced sum $23 > 21$.
    \item \textbf{Max Achievable Sum Pruning}: Branch $[3, 5]$ was pruned because the maximum sum achievable with all remaining elements was only $20 < 21$.
    \item \textbf{Cardinality Underflow Pruning}: Branches with insufficient remaining elements to reach $k_{\min}=2$ were pruned immediately.
\end{itemize}

% ============================================================================
% EXPERIMENT 3
% ============================================================================
\section{Experiment 3: Graph Coloring Using Backtracking (Question 3.1)}

\subsection{Problem Statement \& Requirements}
Given an undirected graph $G = (V, E)$, assign colors to its vertices such that no two adjacent vertices share the same color. The algorithm extends standard $m$-coloring by supporting:
\begin{itemize}
    \item \textbf{Maximum colors bound}: At most $m$ available colors (integers $1, 2, \dots, m$).
    \item \textbf{Allowed color lists}: Each vertex $v$ has a restricted set $\text{allowed\_colors}[v] \subseteq \{1, \dots, m\}$.
    \item \textbf{Pre-colored vertices}: A set of fixed color assignments $\text{pre\_colored}[v]$ that cannot be changed.
    \item \textbf{Adjacency conflict rejection}: Reject assignments where adjacent vertices receive the same color.
    \item \textbf{Constraint pruning}: Prune branches immediately when a color choice violates an adjacency or allowed list constraint.
    \item \textbf{Search statistics}: Report the total number of recursive calls made and branches pruned.
\end{itemize}

\subsection{Algorithmic Formulation}
Vertices are indexed $0, \dots, |V|-1$. The recursive function \texttt{backtrack(idx)} colors vertex $u = V[\text{idx}]$:
\begin{itemize}
    \item If vertex $u$ is pre-colored, verify that its pre-assigned color does not conflict with any colored neighbors in $G$. If safe, recurse to vertex $\text{idx}+1$; otherwise, increment \texttt{pruned\_branches} and backtrack.
    \item If vertex $u$ is uncolored, iterate through colors $c \in \text{allowed\_colors}[u]$ where $c \le m$. If $c$ does not conflict with any adjacent neighbor (\texttt{is\_safe(u, c)}), tentatively assign $\text{color}[u] = c$ and recurse; otherwise, increment \texttt{pruned\_branches}.
\end{itemize}

\subsection{Python Source Code Implementation}
The complete implementation of the list-constrained graph coloring solver is shown below:

\begin{lstlisting}[style=code]
def solve_list_constrained_graph_coloring(adj, m, allowed_colors, pre_colored=None):
    """
    Backtracking solver for List-Constrained Graph Coloring.
    
    Parameters:
    - adj: dict {vertex: [neighbors]} representing an undirected graph.
    - m: maximum available colors (1 to m).
    - allowed_colors: dict {vertex: [allowed_color_integers]}.
    - pre_colored: dict {vertex: fixed_color}.
    
    Returns:
    - solutions: list of valid colorings (each is a dict {vertex: color}).
    - stats: dict containing 'recursive_calls' and 'pruned_branches'.
    """
    if pre_colored is None:
        pre_colored = {}

    vertices = list(adj.keys())
    n = len(vertices)

    # Requirement 4 & Input validation: Check pre-colored vertices
    for u, col in pre_colored.items():
        if col not in allowed_colors.get(u, []):
            raise ValueError(f"Contradictory input: Pre-colored vertex {u} has color {col} not in its allowed list.")
        if col > m or col < 1:
            raise ValueError(f"Contradictory input: Pre-colored vertex {u} has color {col} exceeding available colors {m}.")
        # Check pre-colored adjacency conflict
        for v in adj[u]:
            if v in pre_colored and pre_colored[v] == col:
                raise ValueError(f"Contradictory input: Adjacent pre-colored vertices {u} and {v} both have color {col}.")

    colors = {u: None for u in vertices}
    for u, col in pre_colored.items():
        colors[u] = col

    solutions = []
    stats = {"recursive_calls": 0, "pruned_branches": 0}

    def is_safe(u, col):
        # Requirement 3: Reject any assignment in which adjacent vertices share a color
        for v in adj[u]:
            if colors[v] == col:
                return False
        return True

    def backtrack(idx):
        stats["recursive_calls"] += 1

        # Base case: All vertices colored
        if idx == n:
            solutions.append(dict(colors))
            return

        u = vertices[idx]

        # If vertex is pre-colored, preserve its color and verify adjacency
        if u in pre_colored:
            if is_safe(u, pre_colored[u]):
                backtrack(idx + 1)
            else:
                stats["pruned_branches"] += 1
            return

        # Requirement 5: Try only colors from the vertex's allowed list up to m
        for col in allowed_colors.get(u, []):
            if col > m:
                continue
            if is_safe(u, col):
                colors[u] = col
                backtrack(idx + 1)
                colors[u] = None
            else:
                # Requirement 6: Prune branch as soon as color violates adjacency
                stats["pruned_branches"] += 1

    backtrack(0)
    return solutions, stats\end{lstlisting}

\subsection{Demonstration Test Code}
The demonstration code evaluates a 5-vertex undirected graph with cycle and cross edges, specific vertex allowed lists, and a pre-colored vertex:

\begin{lstlisting}[style=code]
print("=== Experiment 3: List-Constrained Graph Coloring Demonstrations ===")

# Undirected graph G with 5 vertices (V0 to V4):
# Edges: (0-1), (0-2), (1-2), (1-3), (2-4), (3-4)
adj = {
    0: [1, 2],
    1: [0, 2, 3],
    2: [0, 1, 4],
    3: [1, 4],
    4: [2, 3]
}

m = 3  # Colors: 1, 2, 3

# Vertex-specific allowed color lists
allowed_colors = {
    0: [1, 2],
    1: [1, 2, 3],
    2: [2, 3],
    3: [1, 3],
    4: [1, 2, 3]
}

# Pre-colored vertex: Vertex 0 is fixed to color 1
pre_colored = {0: 1}

colorings, stats = solve_list_constrained_graph_coloring(adj, m, allowed_colors, pre_colored)

print(f"Graph vertices: {list(adj.keys())}")
print(f"Max colors (m): {m}")
print(f"Allowed colors per vertex: {allowed_colors}")
print(f"Pre-colored vertices: {pre_colored}")
print("-" * 50)
print(f"Valid Colorings Found: {len(colorings)}")
for i, col_map in enumerate(colorings, 1):
    print(f"  Coloring {i}: {col_map}")

print("-" * 50)
print("Search Statistics:")
print(f"  Total recursive calls: {stats['recursive_calls']}")
print(f"  Total branches pruned: {stats['pruned_branches']}")

# Test Contradictory Input: Adjacent vertices pre-colored with same color
print("\nTesting Contradictory Input Detection:")
try:
    solve_list_constrained_graph_coloring(adj, m, allowed_colors, pre_colored={0: 1, 1: 1})
except ValueError as e:
    print("  Successfully rejected:", e)\end{lstlisting}

\subsection{Execution Results \& Output Logs}
The execution outputs including valid colorings, search statistics, and contradictory input handling are shown below:

\begin{lstlisting}[style=output]
=== Experiment 3: List-Constrained Graph Coloring Demonstrations ===
Graph vertices: [0, 1, 2, 3, 4]
Max colors (m): 3
Allowed colors per vertex: {0: [1, 2], 1: [1, 2, 3], 2: [2, 3], 3: [1, 3], 4: [1, 2, 3]}
Pre-colored vertices: {0: 1}
--------------------------------------------------
Valid Colorings Found: 4
  Coloring 1: {0: 1, 1: 2, 2: 3, 3: 1, 4: 2}
  Coloring 2: {0: 1, 1: 2, 2: 3, 3: 3, 4: 1}
  Coloring 3: {0: 1, 1: 2, 2: 3, 3: 3, 4: 2}
  Coloring 4: {0: 1, 1: 3, 2: 2, 3: 1, 4: 3}
--------------------------------------------------
Search Statistics:
  Total recursive calls: 13
  Total branches pruned: 9

Testing Contradictory Input Detection:
  Successfully rejected: Contradictory input: Adjacent pre-colored vertices 0 and 1 both have color 1.

\end{lstlisting}

\subsection{Analysis}
\begin{itemize}
    \item \textbf{Constrained Search Efficiency}: For the 5-vertex test graph with $m=3$, restricted color domains and pre-colored vertex $0 \mapsto 1$ yielded exactly 4 valid vertex colorings.
    \item \textbf{Search Statistics}: The search required 13 recursive calls and pruned 9 branches as soon as adjacency collisions or color list restrictions occurred.
    \item \textbf{Contradictory Pre-coloring Rejection}: Pre-coloring adjacent vertices $0$ and $1$ with identical color $1$ was caught and rejected during input validation before recursion.
\end{itemize}

% ============================================================================
% CONCLUSION
% ============================================================================
\section{Conclusion}
This lab demonstrated the effectiveness of backtracking in solving complex combinatorial search problems with domain and structural constraints. Across all three experiments (Constrained 8-Queens, Constrained Sum of Subsets, and List-Constrained Graph Coloring):
\begin{enumerate}
    \item \textbf{Early Bounding}: Checking constraints prior to recursive descent dramatically shrinks the active search space compared to brute-force enumeration.
    \item \textbf{Constraint Pruning}: Combining problem-specific bounds (e.g., maximum achievable sum, cardinality limits, diagonal collision checks) ensures early termination of infeasible branches.
    \item \textbf{Input Validation}: Detecting contradictory constraints before initiating search preserves computational resources and avoids invalid states.
\end{enumerate}

\end{document}
